\documentclass[12pt,a4paper]{article}
\usepackage[utf8]{inputenc}
\usepackage[T1]{fontenc}
\usepackage[brazil]{babel}
\usepackage{amsmath}
\usepackage{amssymb}
\usepackage{geometry}
\geometry{margin=2.5cm}
\usepackage{listings}
\usepackage{xcolor}

\lstset{
    language=Python,
    basicstyle=\ttfamily\small,
    keywordstyle=\color{blue},
    stringstyle=\color{red},
    commentstyle=\color{green!60!black},
    numbers=none,
    backgroundcolor=\color{gray!10},
    showspaces=false,
    showstringspaces=false,
    showtabs=false,
    frame=single,
    tabsize=4,
    captionpos=b,
    breaklines=true,
    breakatwhitespace=false,
    escapeinside={\%*}{*)}
}

\begin{document}

\title{Relatório \\ Reconhecimento Facial com Eigenfaces (Implementação via PCA)}
\author{Daniel Rossano Piccoli de Oliveira \\ DRE: 124174238}
\date{}
\maketitle

\noindent\rule{\textwidth}{0.4pt}

\section*{Resumo}

Esse relatório descreve a implementação de um sistema de
reconhecimento facial simplificado baseado na técnica de Eigenfaces (Turk \&
Pentland, 1991), desenvolvida em Python. O objetivo do traballho foi aplicar PCA para comprimir rostos num conjunto pequeno de números e usar essa representação comprimida para identificar pessaoas. O dataset usado foi o \textit{Olivetti Faces} com 40 pessoas e 10 fotos cada (com expressões, ângulos e iluminação ligeiramente diferentes em cada foto) totalizando 400 imagens.
O núcleo matemático que envolveu o cálculo de componentes principais, projeção, reconstrução e
classificação foi implementado a partir dos princípios aprendidos em
CoCaDA sem uso de bibliotecas de alto nível que encapsulem o PCA. Cada seção deste relatório corresponde exatamente a uma seção do
notebook do Colab onde o projeto foi apresentado.

\noindent\rule{\textwidth}{0.4pt}

\section{Carregamento dos dados}

Uma imagem digital em escala de cinza de dimensão $64 \times 64$ é, formalmente, uma
função discreta $I: \{0,\dots,63\}\times\{0,\dots,63\} \to \mathbb{R}$. Para que os
métodos de álgebra linear (autovalores, produto interno e normas) sejam
aplicáveis, é necessário reinterpretar essa estrutura bidimensional como um vetor
em $\mathbb{R}^{4096}$, ou seja, uma operação de \textit{vetorização} (empilhamento das linhas ou
colunas da imagem em sequência), que é uma bijeção linear e não descarta nenhuma
informação, apenas reorganiza os mesmos 4096 valores escalares em uma estrutura de
dado unidimensional.

Reunindo $n=400$ dessas vetorizações como colunas, obtém-se a matriz de dados:

$$A \in \mathbb{R}^{4096 \times 400}, \qquad A = \begin{bmatrix} a_1 & a_2 & \cdots & a_{400}\end{bmatrix}$$

onde cada coluna $a_i \in \mathbb{R}^{4096}$ é um rosto.

\begin{lstlisting}
caminho_npz = "/content/olivetti_faces.npz"
dados_npz = np.load(caminho_npz)
imagens = dados_npz["imagens"] 
rotulos = dados_npz["rotulos"] 

altura, largura = imagens.shape[1], imagens.shape[2]
A = imagens.reshape(imagens.shape[0], -1).T
n_pixels, n_amostras = A.shape
n_pessoas = len(set(rotulos))
\end{lstlisting}

\begin{itemize}
  \item \texttt{np.load(caminho\_npz)} carrega um arquivo binário do Olivetti Faces já subido contendo
  o conjunto de imagens (\texttt{imagens}, array de forma \texttt{(400, 64, 64)}) e seus rótulos de
  identidade (\texttt{rotulos}, array de forma \texttt{(400,)}, com valores inteiros de 0 a 39
  identificando cada uma das 40 pessoas).
  \item \texttt{imagens.reshape(imagens.shape[0], -1)} executa a vetorização descrita
  acima, transformando o array tridimensional \texttt{(400, 64, 64)} em um array bidimensional
  \texttt{(400, 4096)}, onde a dimensão \texttt{-1} é inferida automaticamente pelo NumPy como
  $64 \times 64 = 4096$.
  \item O \texttt{.T} (transposição) subsequente converte essa matriz de \texttt{(400, 4096)} para
  \texttt{(4096, 400)}. Essa transposição garante que cada coluna, e não cada
  linha, represente um rosto.
  \item \texttt{n\_pixels, n\_amostras = A.shape} desempacota as dimensões resultantes (4096 e 400) para uso posterior em todo o notebook.
  \item \texttt{len(set(rotulos))} calcula a o número de rótulos distintos,
  confirmando que existem 40 identidades no conjunto de dados.
\end{itemize}

\section{Rosto médio e centralização}

O Método das Componentes Principais busca as direções ortogonais de máxima variância
de um conjunto de dados. Antes de aplicar o PCA, é preciso centralizar os dados, ou seja, calcular a média de cada pixel ao longo das 400 fotos e obter o "rosto médio" que é um rosto abstrato que representa a média entre todos os rostos. Depois se subtrai essa média de cada elemento da matriz. Isso é necessário porque o PCA busca direções de máxima variância dos dados e a variância é definida em relação à média. Sem centraizar, a direção de maior variação ficaria contaminada pela posição do centro de massa dos dados e não pela variaç\textasciitilde ao real entre os rostos. Com o rosto médio, qualquer variação a partir dele será uma característica individual de cada pessoa.

Define-se o rosto médio por:

$$\bar{a} = \frac{1}{n}\sum_{i=1}^{n} a_i \in \mathbb{R}^{4096}$$

e a matriz de dados centralizada:

$$A_c = A - \bar{a}\,\mathbf{1}_n^T, \quad \text{ou seja,} \quad (A_c)_{:,i} = a_i - \bar{a} \;\; \forall i$$

\begin{lstlisting}
media = A.mean(axis=1)
A_centralizada = A - media[:, None]
\end{lstlisting}

\begin{itemize}
  \item \texttt{A.mean(axis=1)} calcula a média ao longo do eixo 1 (as colunas) para cada uma das
  4096 linhas. Para cada pixel, calcula-se a média de seu valor através das
  400 fotos. O resultado é o vetor $\bar{a}$, de forma \texttt{(4096,)}.
  \item \texttt{media[:, None]} é uma operação de reshape implícita que insere um eixo adicional,
  transformando o vetor de forma \texttt{(4096,)} em uma matriz-coluna de forma \texttt{(4096, 1)}.
  Essa transformação é necessária para que a subtração seguinte funcione
  corretamente. O NumPy expande automaticamente um array de forma
  \texttt{(4096, 1)} para operar elemento a elemento contra um array de forma \texttt{(4096, 400)},
  subtraindo o vetor médio de cada uma das 400 colunas simultaneamente, sem
  necessidade de laços explícitos.
  \item O resultado, \texttt{A\_centralizada}, é a $A_c$ definida acima e é a matriz centralizada
  utilizada em toda a análise de componentes principais subsequente.
\end{itemize}

A célula correspondente no Colab também plota \texttt{media.reshape(altura, largura)}, remontando o vetor médio de
volta à sua forma bidimensional original, mostrando o rosto médio apenas para fins de explanação visual, sem
nenhuma implicação matemática adicional.

\section{PCA via autovalores de A\textsuperscript{T}A}

Seção central do notebook, tanto em complexidade teórica quanto em
importância para a arquitetura de todo o sistema.

\subsection{Problema de custo do PCA}

Componentes principais são definidos como os autovetores da matriz de covariância
amostral dos dados centralizados:

$$C = \frac{1}{n-1} A_c A_c^T \in \mathbb{R}^{4096 \times 4096}$$

É dividida a matriz (n-1) para que os autovalores resultantes representem a variância estatística capturada por cada componente ao invés da soma dos quadrados da disperção dos dados.

Resolver o problema de autovalores/autovetores diretamente sobre $C$ exigiria
diagonalizar uma matriz $4096\times4096$, uma operação de alto custo em relação à dimensão $p=4096$, computacionalmente excessiva nesse
contexto.

\subsection{Redução ao problema equivalente em A\textsuperscript{T}A}

Como o número de amostras $n=400$ é muito menor que a dimensão $p=4096$, define-se a
matriz alternativa muito menor:

$$M = \frac{1}{n-1}A_c^TA_c \in \mathbb{R}^{400\times400}$$

*Os autovetores não-nulos de $A_c^TA_c$ estão em correspondência
bijetiva com os autovetores não-nulos de $A_cA_c^T$, através da aplicação de $A_c$,
preservando o autovalor associado.

\textbf{Demonstração:} seja $v \in \mathbb{R}^{400}$ um autovetor de $A_c^TA_c$ associado
ao autovalor $\lambda$, isto é:

$$A_c^TA_c\,v = \lambda v$$

Multiplicando ambos os lados dessa igualdade por $A_c$:

$$A_c\left(A_c^TA_c\,v\right) = A_c(\lambda v)$$
$$\left(A_cA_c^T\right)(A_cv) = \lambda(A_cv)$$

Essa última expressão afirma
que o vetor $w = A_cv \in \mathbb{R}^{4096}$ satisfaz $A_cA_c^T w = \lambda w$, então $w$ é autovetor de $A_cA_c^T$ associado ao mesmo autovalor $\lambda$.
$\blacksquare$

O vetor $w=A_cv$ obtido pela
proposição acima não possui, em geral, norma unitária, ainda que $v$ tenha sido
normalizado (a aplicação linear $A_c$ não preserva norma). É necessário, portanto,
normalizar explicitamente $w$ após sua obtenção, para que constitua um autovetor no
sentido próprio da palavra (com norma 1), adequado a servir como vetor de uma base
ortonormal.

\begin{lstlisting}
def calcular_pca(A, n_componentes=None):
    n_amostras = A.shape[1]
    if n_componentes is None:
        n_componentes = n_amostras

    M = (A.T @ A) / (n_amostras - 1)
    autovalores, autovetores_M = np.linalg.eigh(M)

    ordem = np.argsort(autovalores)[::-1]
    autovalores = autovalores[ordem]
    autovetores_M = autovetores_M[:, ordem]

    autovetores = A @ autovetores_M
    normas = np.linalg.norm(autovetores, axis=0)
    normas[normas == 0] = 1
    autovetores = autovetores / normas

    n_componentes = min(n_componentes, autovetores.shape[1])
    return autovetores[:, :n_componentes], autovalores[:n_componentes]
\end{lstlisting}

\begin{itemize}
  \item \texttt{M = (A.T @ A) / (n\_amostras - 1)}: constrói explicitamente a matriz $M$ de dimensão $400\times400$. O parâmetro \texttt{A} dessa função
  recebe, na prática, a matriz \texttt{A\_centralizada} já centralizada.
  \item \texttt{np.linalg.eigh(M)} retorna uma tupla \texttt{(autovalores, autovetores\_M)}. Por convenção
  da biblioteca, \texttt{eigh} devolve os autovalores em ordem \textbf{crescente}; \texttt{autovetores\_M}
  é uma matriz cujas colunas são os autovetores correspondentes de $M$.
  \item \texttt{ordem = np.argsort(autovalores)[::-1]}: \texttt{np.argsort} produz os índices que
  ordenariam o array de forma crescente, a fatia \texttt{[::-1]} inverte essa sequência de
  índices, produzindo a ordem decrescente desejada.
  \item \texttt{autovetores = A @ autovetores\_M}: essa linha implementa a operação
  $w=A_cv$. Para cada coluna $v$ de
  \texttt{autovetores\_M}, calcula-se o produto matriz-vetor correspondente, recuperando os
  400 autovetores (não-normalizados) de $A_cA_c^T$, cada um de dimensão 4096.
  \item \texttt{normas = np.linalg.norm(autovetores, axis=0)}: calcula a norma de cada
  coluna independentemente (\texttt{axis=0} percorre as linhas dentro de cada coluna fixa).
  A linha seguinte, \texttt{normas[normas == 0] = 1}, é uma proteção contra
  divisão por zero, para o caso teoricamente possível (mas imporvável na prática com esse conjunto de dados).
  \item \texttt{autovetores = autovetores / normas}: normalização final, dividindo cada coluna
  pela sua respectiva norma, produzindo autovetores de
  norma unitária (as \textit{eigenfaces}).
  \item As duas últimas linhas realizam a truncagem para os \texttt{n\_componentes} solicitados. Se
  \texttt{n\_componentes=None} foi fornecido, \texttt{n\_componentes} foi redefinido como
  \texttt{n\_amostras} (400) no início da função e a fatia final devolve, na prática, todos
  os componentes disponíveis (o \texttt{min} com \texttt{autovetores.shape[1]} é uma garantia
  adicional).
\end{itemize}

No fim, a célula faz

\begin{lstlisting}
autovetores_completos, autovalores_completos = calcular_pca(A_centralizada, n_componentes=None)
\end{lstlisting}

obtendo a base completa de componentes principais que, na prática, tem
exatamente 399 componentes não-triviais. A operação de centralização tem
$\sum_i (A_c)_{:,i} = \mathbf{0}$ (a soma das colunas
centralizadas é nula, pela média subtraída), o que
reduz o posto máximo de $A_c$ em uma unidade em relação ao número de
colunas, de $400$ para $399$.

\section{Visualização das eigenfaces}

As eigenfaces são os autovetores finais da matriz de covariância remontados como imagens 64 por 64 revelando as direções de variações entre os rostos (como presença de óculos, iluminação, formato geral do rosto, etc.). As eigenfaces não são rostos reais, são apenas direções que capturam o máximo de variância estatística possível.

\begin{lstlisting}
fig, eixos = plt.subplots(2, 5, figsize=(14, 6))
for i, eixo in enumerate(eixos.flat):
    eixo.imshow(autovetores_completos[:, i].reshape(altura, largura), cmap="gray")
    eixo.set_title(f"Eigenface {i+1}")
    eixo.axis("off")
\end{lstlisting}

Cada \texttt{autovetores\_completos[:, i]} é um vetor de $\mathbb{R}^{4096}$ com norma 1. O
\texttt{.reshape(altura, largura)} desfaz a vetorização descrita, remontando o
vetor em sua forma matricial bidimensional original $64\times64$ para
fins de visualização das imagens.

Cada eigenface representa uma direção
ortogonal no espaço $\mathbb{R}^{4096}$ ao longo da qual os rostos do conjunto de
treinamento exibem variação estatística decrescente (a primeira eigenface,
correspondente ao maior autovalor, é a direção de maior variância e assim vai).

\section{Variância explicada}

CAda autovalor associado a uma eigenface representa quanta variância dos dados originais está concentrada naquela direção. Somando os autovalores dos primeiros k componentes e dividindo pela soma total, obtém-se a variância explicada acumulada. Essa métrica responde a pergunta de "quantos componentes é preciso manter para não perder muita informação". No experimento aplicado no código, cerca de 66 componentes já capturam 90\% da variância total, 123 capturam 95\% e 260 capturam 99\% de um total de 399 componentes possíveis.

\begin{lstlisting}
def variancia_explicada_acumulada(autovalores):
    autovalores = np.clip(autovalores, a_min=0, a_max=None)
    total = autovalores.sum()
    return np.cumsum(autovalores) / total if total > 0 else np.zeros_like(autovalores)
\end{lstlisting}

\begin{itemize}
  \item \texttt{np.clip(autovalores, a\_min=0, a\_max=None)}: impõe um piso de zero aos autovalores.
  Teoricamente, todos os autovalores de uma matriz positiva (como $M$)
  são não-negativos. Essa linha é uma proteção contra autovalores ligeiramente
  negativos que podem surgir por erro de arredondamento de ponto flutuante em
  componentes cujo autovalor teórico é zero.
  \item \texttt{np.cumsum(autovalores)}: calcula a soma cumulativa até o componente de posição k.
  \item A divisão pelo \texttt{total} produz o vetor completo para todo $k$ simultaneamente.
\end{itemize}

\begin{lstlisting}
def n_componentes_para_variancia(autovalores, alvo):
    acumulada = variancia_explicada_acumulada(autovalores)
    indices = np.where(acumulada >= alvo)[0]
    return int(indices[0]) + 1 if len(indices) else len(autovalores)
\end{lstlisting}

Essa função resolve o problema inverso, ou seja, encontrar o menor $k$ tal
que a variância explicada seja maior ou igual ao alvo \texttt{np.where(acumulada >= alvo)[0]} retorna todos
os índices que satisfazem a condição. Como o vetor \texttt{acumulada} é não-decrescente por construção (soma cumulativa de valores
não-negativos), o primeiro desses índices, somado a 1 (para corrigir a 0-indexação
do NumPy para a contagem de número de componentes) é a resposta.

\section{Reconstrução com rank crescente}

Dada a base ortonormal completa de eigenfaces $\{v_1,\dots,v_{399}\}$, qualquer rosto
centralizado $a-\bar a$ admite a representação exata:

$$a - \bar a = \sum_{i=1}^{399} c_i v_i, \qquad c_i = v_i^T(a-\bar a)$$

por se tratar de uma expansão em base ortonormal completa do subespaço gerado pelos
dados. A \textbf{aproximação de posto $k$} consiste em truncar esta soma nos $k$ primeiros
termos:

$$\hat a_k = \bar a + \sum_{i=1}^{k} c_i v_i = \bar a + V_k V_k^T(a-\bar a)$$

onde $V_k = \begin{bmatrix}v_1&\cdots&v_k\end{bmatrix}$.

Um rosto pode ser aproximado usando só os k primeiros componentes. Projetar significa trasformar um rosto de 4096 pixels em um vetor pequeno de k coeficientes, multiplicando pela transposta das eigenfaces. Reconstruir é o caminho inverso, multiplicar os coeficientes pelas eigenfaces e somar de volta a média, retornando um rosto aproximado de acordo com a quantidade de componentes usados.

\subsection*{O que o código faz}

As funções \texttt{projetar} e \texttt{reconstruir}, definidas na Seção 3 do notebook, implementam
exatamente as duas metades desta equação:

\begin{lstlisting}
def projetar(A, media, autovetores):
    A_centralizada = A - media[:, None]
    return autovetores.T @ A_centralizada

def reconstruir(coeficientes, media, autovetores):
    return autovetores @ coeficientes + media[:, None]
\end{lstlisting}

\texttt{projetar} calcula $c = V_k^T(a-\bar a)$: centraliza a entrada subtraindo a média
recebida como parâmetro, e multiplica pela transposta da matriz de eigenfaces
fornecida (\texttt{autovetores} $(4096,k)$), resultando em um vetor de
coeficientes de forma $(k,)$ que é a representação comprimida do rosto.

\texttt{reconstruir} calcula $\hat a = V_kc+\bar a$, multiplicando os coeficientes pela matriz
de eigenfaces (sem transpor, retornando ao espaço de 4096 dimensões) e soma de volta
a média, desfazendo a centralização.

A função \texttt{projetar} já realiza a
centralização internamente, subtraindo \texttt{media} de seu argumento de entrada.
Consequentemente, todas as chamadas a \texttt{projetar} no notebook recebem como
argumento os dados originais, não centralizados (por exemplo, \texttt{A} e não
\texttt{A\_centralizada}). Fornecer dados já centralizados resultaria em subtração da média
duas vezes, distorcendo a representação comprimida resultante.

\begin{lstlisting}
ranks = [1, 5, 50, 150, 299]
for r in ranks:
    coef = projetar(A[:, [indice_exemplo]], media, autovetores_completos[:, :r])
    recon = reconstruir(coef, media, autovetores_completos[:, :r])
\end{lstlisting}

\texttt{autovetores\_completos[:, :r]} é a operação que concretiza $V_r$ que é a
matriz contendo apenas as primeiras $r$ colunas (eigenfaces) da base completa. É
neste ponto que a redução de dimensionalidade ocorre. Ao aumentar $r$ progressivamente (1, 5, 50, 150, 299), observa-se
visualmente a reconstrução tornar-se mais fiel ao rosto original, com
o rank máximo (299) produzindo reconstrução visualmente quase idêntica à
original ( o que já era esperado, pois $r=299$ já captura a maioria da
variância total (superior a 99\%), estando muito próximo do posto
máximo $n-1=399$ que conteria a totalidade da informação representável pela base.

\section{Erro de reconstrução vs. número de componentes}

O Teorema de Eckaert-Young-Mirsky estabelece que, para qualquer
matriz $A_c$ e qualquer $k \leq \text{posto}(A_c)$, a aproximação de posto $k$ que
minimiza o erro na norma de Frobenius,

$$\min_{\text{posto}(B)\leq k} \|A_c - B\|_F$$

é $B_k^* = V_kV_k^TA_c$ que é a projeção de $A_c$ sobre o subespaço gerado
pelos $k$ primeiros componentes principais. Uma consequência desse teorema é
que a função de erro é não crescente, então aumentar o número de componentes nunca vai piorar a
melhor aproximação possível, apenas mantê-la ou melhorá-la.

\begin{lstlisting}
erros = []
lista_ranks_erro = list(range(1, 51, 2))
for r in lista_ranks_erro:
    coef = projetar(A, media, autovetores_completos[:, :r])
    recon = reconstruir(coef, media, autovetores_completos[:, :r])
    erros.append(np.linalg.norm(A - recon, ord="fro"))
\end{lstlisting}

Para cada valor de $r$ na sequência $\{1,3,5,\dots,49\}$, a *matriz
inteira* de 400 rostos é projetada e reconstruída e o erro é medido pela norma de Frobenius da diferença entre a matriz
original e a reconstruída:

$$\|A - \hat A_r\|_F = \sqrt{\sum_{i,j}(A_{ij}-\hat A_{r,ij})^2}$$

\begin{lstlisting}
decrescente = all(erros[i] >= erros[i+1] - 1e-6 for i in range(len(erros) - 1))
print("Erro progressivamente decrescente?", decrescente)
\end{lstlisting}

Essa linha realiza uma verificação computacional da propriedade
teórica do Teorema de Eckart-Young: testa, para cada par de valores
de rank consecutivos na sequência testada, se o erro não aumentou (com tolerância de
$10^{-6}$). O resultado \texttt{True} obtido na
execução prova que a implementação está
correta.

\section{Classificação por vizinho mais próximo}

O reconhecimento é resovido como um problema de classificação de padrões pelo método do vizinho mais próximo: dado um vetor de consulta $c_{\text{teste}}$ no espaço
reduzido, atribui-se a ele o rótulo do vetor de treinamento $c_{\text{treino}}^{(j)}$
que minimiza a distância euclidiana:

$$j^* = \arg\min_j \left\| c_{\text{treino}}^{(j)} - c_{\text{teste}} \right\|_2$$

A avaliação de um classificador exige que o desempenho seja medido sobre dados que não participaram de
seu ajuste. Se a base de eigenfaces fosse calculada utilizando todas as 400
fotos (incluindo aquelas escolhidas para teste), a estimativa de desempenho estaria
contaminada por vazamento de dados, a base de componentes
principais teria incluído as fotos de teste, aumentando artificialmente e invalidando a precisão relatada.

\begin{lstlisting}
def classificar_vizinho_mais_proximo(coef_treino, rotulos_treino, coef_teste):
    previsoes = []
    for i in range(coef_teste.shape[1]):
        vetor_teste = coef_teste[:, i]
        distancias = np.linalg.norm(coef_treino - vetor_teste[:, None], axis=0)
        previsoes.append(rotulos_treino[np.argmin(distancias)])
    return np.array(previsoes)
\end{lstlisting}

Para cada coluna \texttt{i} de \texttt{coef\_teste} (cada foto de teste projetada no espaço
reduzido), calcula-se \texttt{coef\_treino - vetor\_teste[:, None]}: 
subtrai-se o vetor de teste (forma \texttt{(k,1)} após o \texttt{[:, None]}) de cada coluna de
\texttt{coef\_treino} (forma \texttt{(k, n\_treino)}) simultaneamente, produzindo uma matriz de
diferenças de mesma forma que \texttt{coef\_treino}. \texttt{np.linalg.norm(..., axis=0)} calcula a
norma euclidiana de cada coluna dessa matriz de diferenças independentemente,
produzindo o vetor \texttt{distancias}, de forma \texttt{(n\_treino,)}, contendo a distância da foto
de teste a cada uma das fotos de treino. \texttt{np.argmin(distancias)} localiza o
índice de menor distância e o id da foto escolhida é anexado à lista de previsões.

\begin{lstlisting}
idx_treino, idx_teste = train_test_split(
    indices, test_size=0.25, stratify=rotulos, random_state=42
)
\end{lstlisting}

\texttt{train\_test\_split} separa os 400 índices em dois subconjuntos disjuntos:
300 (75\%) para treino e 100 (25\%) para teste. O parâmetro
\texttt{stratify=rotulos} garante que a proporção de cada uma das 40 classes seja preservada
nos dois subconjuntos. \texttt{random\_state=42} fixa a semente do
gerador pseudoaleatório, permitindo que a divisão entre
execuções distintas do notebook seja reproduzível.

\begin{lstlisting}
media_treino = A_treino.mean(axis=1)
A_treino_centralizada = A_treino - media_treino[:, None]

n_max = min(A_treino.shape[1] - 1, 150)
autovetores_treino, _ = calcular_pca(A_treino_centralizada, n_componentes=n_max)
\end{lstlisting}

A média (\texttt{media\_treino}) e a base de eigenfaces (\texttt{autovetores\_treino}) são calculadas
exclusivamente a partir de \texttt{A\_treino}, nenhuma informação de \texttt{A\_teste}
participa desse cálculo. \texttt{n\_max = min(A\_treino.shape[1] - 1, 150)} limita o número de componentes
calculados a no máximo 150 (por razões de custo computacional. O posto máximo de
\texttt{A\_treino\_centralizada} seria de 299,
mas apenas 150 são efetivamente necessários dado o comportamento de saturação já
antecipado pela análise de variância explicada).

\begin{lstlisting}
n_componentes_clf = 50
coef_treino = projetar(A_treino, media_treino, autovetores_treino[:, :n_componentes_clf])
coef_teste = projetar(A_teste, media_treino, autovetores_treino[:, :n_componentes_clf])
\end{lstlisting}

Ambas as chamadas de \texttt{projetar} usam \texttt{media\_treino} e a mesma fatia de
\texttt{autovetores\_treino} que reutiliza os
parâmetros aprendidos no treino sem jamais recalcular estatística nenhuma a partir do
próprio conjunto de teste. Essa é a materialização em código da restrição:
os "parâmetros aprendidos" do sistema são exatamente o par
$(\bar a_{\text{treino}}, V_{50})$, e o teste consiste em avaliar se a projeção de
fotos não vistas durante este cálculo, usando esses parâmetros fixos, ainda preserva
separabilidade suficiente entre as 40 classes para uma classificação por distância
correta.

\begin{lstlisting}
previsoes = classificar_vizinho_mais_proximo(coef_treino, rotulos_treino, coef_teste)
precisao = np.mean(previsoes == rotulos_teste)
\end{lstlisting}

\texttt{previsoes == rotulos\_teste} produz um array booleano de comparação elemento a
elemento. \texttt{np.mean} sobre um array booleano calcula a fração de valores \texttt{True}, resultando
exatamente na taxa de acerto, a precisão da classificação. O resultado obtido nesta
execução, com $k=50$ componentes, foi de \textbf{94\%}.

\section{Precisão vs. número de componentes}

Essa seção generaliza o experimento da seção anterior, testando o efeito da
dimensão do espaço reduzido sobre a taxa de acerto:

\begin{lstlisting}
lista_ranks_clf = sorted(set(r for r in [1, 2, 5, 10, 20, 30, 50, 75, 100, n_max] if r <= n_max))
precisoes = []

for r in lista_ranks_clf:
    coef_tr = projetar(A_treino, media_treino, autovetores_treino[:, :r])
    coef_te = projetar(A_teste, media_treino, autovetores_treino[:, :r])
    preds = classificar_vizinho_mais_proximo(coef_tr, rotulos_treino, coef_te)
    precisoes.append(np.mean(preds == rotulos_teste))
\end{lstlisting}

Aqui o custo de recomputar o PCA para cada valor de $r$ testado é evitado: a
base completa \texttt{autovetores\_treino} (calculada uma única vez até 150
componentes) é refatiada (\texttt{autovetores\_treino[:, :r]}) para cada valor
de $r$ na lista de ranks a testar, reaproveitando o resultado da decomposição
espectral já obtida, e testando exclusivamente o efeito da redução de
dimensionalidade sobre a classificação, sem custo adicional de re-diagonalização.

A precisão cresce rapidamente com poucos componentes e
satura em torno de 20 a 30 componentes (evidência empírica de que a informação
discriminativa entre identidades está concentrada nos componentes principais de
maior variância), consistente com a análise de variância explicada e
sugerindo que os componentes de ordem mais alta (associados a autovalores pequenos)
contribuem predominantemente com ruído estatístico irrelevante para a tarefa de
classificação, não com sinal discriminativo adicional.

\section{Demonstração concreta de uma classificação}

Aqui mostra-se o teste final da classificação, deixando explícito o mecanismo interno de decisão executado de forma agregada sobre as 100 fotos de teste simultaneamente.

\begin{lstlisting}
def testar_uma_classificacao(indice_teste, A_treino, A_teste, media_treino, autovetores,
                              rotulos_treino, rotulos_teste, altura, largura, top_n=5):
    vetor_teste = projetar(A_teste[:, [indice_teste]], media_treino, autovetores)[:, 0]
    coef_treino_local = projetar(A_treino, media_treino, autovetores)

    distancias = np.linalg.norm(coef_treino_local - vetor_teste[:, None], axis=0)
    ordem = np.argsort(distancias)
\end{lstlisting}

\texttt{A\_teste[:, [indice\_teste]]} seleciona uma única coluna de \texttt{A\_teste} (o uso de uma
lista de um único elemento, \texttt{[indice\_teste]}, em vez do inteiro \texttt{indice\_teste}
diretamente, preserva a dimensão de coluna, mantendo o resultado como uma matriz de
forma \texttt{(4096, 1)} ao invés de reduzi-lo a um vetor unidimensional, necessário para
compatibilidade com a assinatura de \texttt{projetar} que espera uma matriz). O restante da
função reaplica a lógica de \texttt{classificar\_vizinho\_mais\_proximo}, mas
retendo, através de \texttt{np.argsort(distancias)}, toda a ordenação das distâncias, e
não apenas o índice de mínimo, permitindo reportar também
as \texttt{top\_n} fotos de treinamento mais próximas, com suas respectivas distâncias
numéricas e uma visualização lado a lado que evidencia visualmente a
plausibilidade (ou implausibilidade) da decisão do classificador. O parâmetro
\texttt{autovetores} recebido por essa função, na chamada realizada na célula
correspondente, é \texttt{autovetores\_treino[:, :n\_componentes\_clf]}, ou seja, a mesma
base de 50 eigenfaces utilizada na classificação agregada, garantindo
que a demonstração individual seja fielmente representativa do experimento
estatístico completo.

\section{Conclusões}

A implementação demonstra que:

\begin{enumerate}
  \item A matriz de covariância de um conjunto de dados de alta dimensão pode ser
  diagonalizada de forma computacionalmente eficiente através da matriz reduzida
  $A^TA$, quando o número de amostras é significativamente menor que a
  dimensionalidade dos dados (resultado justificado pela relação de
  autovalores/autovetores entre $A^TA$ e $AA^T$).
  \item A normalização da covariância por $(n-1)$, e não por $n$, decorre da correção de
  Bessel para estimadores não-viesados de variância, sem afetar a geometria das
  eigenfaces em si.
  \item A aproximação de posto $k$ obtida por truncagem da base de componentes principais
  é, pelo Teorema de Eckart-Young-Mirsky, a melhor aproximação possível na norma de
  Frobenius.
  \item Uma redução de dimensionalidade de 4096 pixels para 50 coeficientes preserva
  informação suficiente para uma tarefa de classificação por vizinho mais próximo
  com 94\% de precisão.
\end{enumerate}

\end{document}